\begin{abstract}
The proliferation of sophisticated cyber-attacks against modern network infrastructures necessitates the deployment of highly accurate Intrusion Detection Systems (IDS). However, real-world network traffic is inherently skewed, characterized by an extreme data imbalance where benign traffic massively overwhelms rare, yet catastrophic, minority attacks. Conventional, solitary machine learning models struggle under these volatile conditions, frequently exhibiting high False Positive (FP) rates while severely compromising the recall of stealthy intrusions. To overcome these critical limitations, this paper proposes a novel, highly optimized framework that synergizes advanced balanced sampling techniques with a sophisticated hybrid ensemble learning architecture. Evaluated on the benchmark CICIDS2017 dataset, our methodology first employs a rigorous hybrid sampling strategy to reconstruct a balanced training vector space, thereby preventing model bias toward majority classes. Subsequently, we introduce a hierarchical ensemble model that leverages the rapid, robust, tree-based gradient boosting capabilities of XGBoost and LightGBM as base learners. The predictive probabilities of these base models are then fed into a Multi-Layer Perceptron (MLP) meta-classifier, which synthesizes non-linear, complex attack patterns. Experimental results demonstrate that the proposed hybrid architecture significantly outperforms state-of-the-art solitary models. By successfully mitigating data imbalance and exploiting the synergistic strengths of gradient boosting and deep learning, the framework achieves a drastic reduction in False Positives while maintaining an exceptionally high recall for critical minority attack classifications. This research contributes a scalable, highly optimized IDS solution capable of rigorous threat detection in realistic, imbalanced network environments.
\end{abstract}
